\section{FORMAL RESTATEMENT}
\textbf{Erd\H{o}s \#379; reconstruction of an explicit prime-power construction, 2026-09-24.}
Let $S(n)$ be the largest integer $r$ such that every internal binomial coefficient $\binom nk$, $1\leq k\leq n-1$, is divisible by an $r$th power of a prime. Is $\limsup_{n\to\infty}S(n)=\infty$?

\section{QUICK LITERATURE AND CONTEXT CHECK}
The current page and discussion record the affirmative elementary construction associated with Cambie, Kova\v c and Tao. The prime may depend on $k$. The argument below shows that just two primes, $2$ and a chosen odd prime $p$, suffice for each constructed row.

\section{OBJECTIVE AND ATTACK PLAN}
Choose the row length as a power of $2$ that is congruent to $1$ modulo $p^r$. One binomial identity handles most indices using their $2$-adic valuation; a second identity handles all remaining indices using $p$.

\section{WORK}
Fix $r\geq2$, choose a prime $p>2^{r-1}$, and set
\[
E=\varphi(p^r)=p^{r-1}(p-1),\qquad n=2^E.
\]
In particular $E\geq r$, and Euler's theorem gives $n\equiv1\pmod {p^r}$. Let $1\leq k<n$.

If $v_2(k)\leq E-r$, the identity
\[
k\binom nk=n\binom{n-1}{k-1}
\]
implies $v_2\binom nk\geq E-v_2(k)\geq r$, so $2^r\mid\binom nk$.

Otherwise $2^{E-r+1}\mid k$. Writing $D=2^{E-r+1}$, we have
\[
k=Dt,\qquad n-k=D(2^{r-1}-t),\qquad1\leq t\leq2^{r-1}-1<p.
\]
Neither $k$ nor $n-k$ is divisible by $p$: both parenthesized factors are positive and less than $p$, and $D$ is a power of $2$. Furthermore $k-1\equiv-(n-k)\pmod p$, so $p\nmid k-1$. Since $k\geq D\geq2$, we can use
\[
k(k-1)\binom nk=n(n-1)\binom{n-2}{k-2}.
\]
The right side is divisible by $p^r$, while $k(k-1)$ is coprime to $p$. Consequently $p^r\mid\binom nk$.

These cases cover every internal index, and prove
\[
S\!\left(2^{\varphi(p^r)}\right)\geq r.                         \tag{1}
\]
The row lengths tend to infinity as $r\to\infty$, proving the required infinite limsup.

One may make the scale explicit using Bertrand's postulate: choose $2^{r-1}<p<2^r$. Then $E=p^{r-1}(p-1)<2^{r^2}$, so for the associated rows
\[
S(n)\geq r>\sqrt{\log_2\log_2 n}.
\]
This last estimate is a subsequence lower bound, not a bound for every row.

\section{RESULT AND LIMITATIONS}
\textbf{Attempt status: solved.} Formula (1) is a complete explicit construction. Euler's theorem and, only for the optional quantitative estimate, Bertrand's postulate are the standard inputs. The proof reconstructs the known solution and makes no novelty claim.

\section{SOURCES}
\url{https://www.erdosproblems.com/379}; the construction and its attribution: \url{https://www.erdosproblems.com/forum/discuss/379}.

\textbf{Completion rate estimate: 100\%.}
